\documentclass[10pt,a4paper]{amsart}
\usepackage[margin=19mm]{geometry}
\usepackage{amsmath,amssymb,mathtools}
\usepackage[scr=rsfs,cal=euler]{mathalfa}
\usepackage{xfrac}
\usepackage[unicode,hidelinks,bookmarks=false]{hyperref}
\newcommand{\ie}{i.e.\ }
\newcommand{\cf}{cf.\ }
\newcommand{\resp}{respectively\ }
\newcommand{\Subsection}{\subsection}
\providecommand{\nobreakdash}{\nobreak\hbox{-}}
\setlength{\emergencystretch}{2em}
\multlinegap0pt
\usepackage{tikz}
\usetikzlibrary{cd}
\usepackage[shortlabels]{enumitem}
\usepackage{mathtools}
\usepackage{amssymb}
\usepackage{xcolor}
\usepackage{float}
\newtheorem{lemma}{Lemma}
\newcommand{\IFF}{\Leftrightarrow}						% Iff
\newcommand{\IMP}{\Rightarrow}						% Left-to-right implication
\newcommand{\Sat}{{\sf Sat}}								%Category of saturated objects
\newcommand{\X}{\ar{X}}									% Class of arrows
\newcommand{\ar}[1]{\mathcal{#1}}							% Class of morphisms
\newcommand{\cat}{{\sf C}}								%A generic category 
\newcommand{\ort}[1]{{#1}^{\perp}}						% Orthogonality class
\title{Beth companions of finitary essentially algebraic theories}
\author{Ivan Di Liberti, Luca Reggio}
\date{Source: arXiv:2609.03601v1 (CC BY 4.0)}
\begin{document}
\maketitle

\noindent\emph{Source and licence.} Beth companions of finitary essentially algebraic theories. Version arXiv:2609.03601v1 (CC BY 4.0), distributed by arXiv under the Creative Commons Attribution 4.0 International licence, http://creativecommons.org/licenses/by/4.0/, as recorded in the arXiv licence metadata for this exact version and shown on the abstract page https://arxiv.org/abs/2609.03601v1. Full licence notice: \texttt{licenses/arxiv-2609.03601-CC-BY-4.0.txt}.\par\smallskip

\noindent\textit{Editorial scope.} Counted unit: the article's lemma l:ort-vs-sat (source0.tex lines 708--715). The article's complete original proof of it is delivered with it (source0.tex lines 717--735, one \texttt{proof} environment of 19 lines). No mathematics is written, repaired or re-derived: the statement and the proof are the article's own lines, in the article's own notation, and no retained line is substituted or re-flowed.  Retained uncounted as the source-local context the proof needs: lines 483--485; lines 646--648.  Nothing was omitted from any retained span, so the package carries no omission record.  No retained line contains a citation, so the wrapper carries no bibliography and no bibliography entry.  The retained lines contain the article's own diagram code, which the wrapper typesets with the standard tikz packages; no image file is delivered and the package declares no asset dependency.  Editorial formatting only, in addition: an AMS \texttt{amsart} preamble that restates the article's own declarations and macros used by the retained lines, namely the environments \texttt{lemma} on separate counters, together with the standard packages those lines need.  The retained environments print in this document as Lemma 1 in order of appearance, whereas the article prints the same items with the numbering of its own full document.  The complete native source of the article as distributed in the arXiv e-print for this exact version is bundled unchanged as \texttt{sources/209308/source0.tex}.
\begin{lemma}\label{l:existence-saturated-epi-ext}
In an lfp category, every object admits a saturated epi-extension.
\end{lemma}

\begin{lemma}\label{l:ort-included-in-sigma}
For any category $\cat$, we have $\ort{\X}\subseteq \Sat(\cat)$.
\end{lemma}

\begin{lemma}\label{l:ort-vs-sat}
The following statements are equivalent for any lfp category $\cat$:
\begin{enumerate}[label=(\arabic*)]
\item\label{i:X-local-extension} every object of $\cat$ admits an extension with codomain in $\ort{\X}$;
\item\label{i:sat-and-ort-coincide} $\Sat(\cat)=\ort{\X}$;
\item\label{i:transferable-epi-monos} $\cat$ has transferable epi-monos.
\end{enumerate}
\end{lemma}

\begin{proof}
\ref{i:X-local-extension} $\IFF$ \ref{i:sat-and-ort-coincide}. As every object of an lfp category has a saturated (epi-)extension by Lemma~\ref{l:existence-saturated-epi-ext}, the implication \ref{i:sat-and-ort-coincide} $\IMP$ \ref{i:X-local-extension} is immediate. For the converse implication, suppose that every object of $\cat$ admits an extension with codomain in $\ort{\X}$. Then $\ort{\X}$ is a mono-reflective subcategory of $\cat$. The unit of this mono-reflection is component-wise epi-monic, hence its component at any saturated object is an isomorphism. It follows that $\Sat(\cat)\subseteq \ort{\X}$, and so $\Sat(\cat)=\ort{\X}$ by Lemma~\ref{l:ort-included-in-sigma}.

\ref{i:sat-and-ort-coincide} $\IFF$ \ref{i:transferable-epi-monos}.
Suppose that $\ort{\X} = \Sat(\cat)$. We claim that $\cat$ has transferable epi-monos. Consider an epi-mono $m\colon Y\to W$ and an arrow $f\colon Y\to X$. Since $\cat$ is lfp, $X$ admits a saturated epi-extension $n\colon X\to Z$. As $\Sat(\cat)\subseteq \ort{\X}$, there is an arrow $g$ making the following diagram commute. Thus, $\cat$ has transferable epi-monos.
\[\begin{tikzcd}
Y \arrow{d}[swap]{m} \arrow{r}{f} & X \arrow{d}{n} \\
W \arrow[dashed]{r}{g} & Z
\end{tikzcd}\]

Conversely, we claim that if $\cat$ has transferable epi-monos, then $\ort{\X} = \Sat(\cat)$ (for this, we do not need to assume that $\cat$ is lfp). 
%
Suppose that $X$ is a saturated object of $\cat$, and consider a span $f,m$ as displayed below, with $m$ epi-mono.
\[\begin{tikzcd}
Y \arrow{d}[swap]{m} \arrow{r}{f} & X \arrow[dashed]{d}{n} \\
W \arrow[dashed]{r}{h} & Z
\end{tikzcd}\]
Since $\cat$ has transferable epi-monos, there exist arrows $h,n$, with $n$ epi-mono, making the above square commute. As $X$ is saturated, $n$ is an isomorphism. Hence, $g\coloneqq n^{-1}\circ h\colon W\to X$ satisfies $g\circ m = f$, showing that $X$ is orthogonal to~$m$. This proves that $\Sat(\cat)\subseteq \ort{\X}$; the other inclusion follows from Lemma~\ref{l:ort-included-in-sigma}.
\end{proof}

\end{document}
